\chapter{Introdução}
\label{capitulo:introducao}

O xadrez é um dos jogos de tabuleiro mais antigos e estudados do mundo, sendo frequentemente utilizado ao longo da história como um laboratório para testar a capacidade lógica e o poder de processamento dos computadores. Nos primórdios da ciência da computação, as partidas contra a máquina eram restritas a terminais de texto escuros ou interfaces extremamente rudimentares. No entanto, com a evolução da tecnologia e a mudança nas expectativas dos jogadores, os jogos digitais passaram a exigir não apenas uma matemática impecável, mas também uma apresentação visual imersiva e responsiva.

Hoje em dia, jogar xadrez contra o computador é algo comum, mas o desenvolvimento desse tipo de software envolve um desafio de arquitetura bastante específico: unir uma interface visual atraente, rica em animações e menus, com um "cérebro" lógico capaz de calcular milhões de probabilidades matemáticas em poucos segundos. 

Motores de jogos modernos, como a Godot Engine, são ferramentas excelentes e acessíveis para criar telas interativas, gerenciar efeitos sonoros e construir menus complexos. Contudo, eles não foram construídos para processar algoritmos pesados de inteligência artificial de forma nativa. Por outro lado, ferramentas consagradas de Inteligência Artificial voltadas para o xadrez, como o Stockfish, são extremamente poderosas e precisas, mas funcionam apenas nos bastidores, através de linhas de comando, sem oferecer nenhuma tela, gráfico ou botão para o jogador humano.

Diante desse cenário, este trabalho apresenta o desenvolvimento do simulador digital \textit{Royal Legacy}. O problema de pesquisa que norteia este projeto é: como integrar de forma eficiente um motor lógico de xadrez de alta complexidade (Stockfish) a uma interface gráfica moderna (Godot Engine), garantindo que o jogo seja visualmente personalizável e que os cálculos da máquina não congelem ou travem a experiência do usuário durante as partidas?

\section{Objetivo Geral}

Este trabalho tem como objetivo desenvolver um simulador de xadrez digital interativo e customizável na Godot Engine, integrado ao motor de inteligência artificial Stockfish através de uma comunicação em linguagem Python, proporcionando uma experiência de jogo fluida, modular e com diferentes níveis de dificuldade.

\section{Objetivos específicos}

Para alcançar o objetivo geral, foram definidos os seguintes objetivos específicos: 

\begin{itemize}
	\item Implementar a interface gráfica e a lógica completa das regras do xadrez (incluindo movimentação restrita, roque e \textit{en passant}) utilizando a linguagem nativa GDScript;
	\item Desenvolver um \textit{middleware} em Python para estabelecer a comunicação assíncrona entre o frontend do jogo e o binário do Stockfish, operando através do protocolo \textit{Universal Chess Interface} (UCI);
	\item Estruturar um sistema de gerenciamento de cenas e menus de opções, garantindo uma experiência de usuário (UX) intuitiva para a configuração de volume, dificuldades e modos de jogo;
	\item Criar um sistema dinâmico de temas visuais baseado em dicionários de dados, permitindo a alteração das cores e dos conjuntos de peças em tempo real;
	\item Desenvolver um gerenciador de estados para controlar turnos e condições de fim de jogo (xeque-mate e empate), validando o fluxo correto das partidas nas modalidades Jogador contra Jogador (PvP) e Jogador contra Máquina (PvE).
\end{itemize}

\section{Justificativa}
\label{secao:justificativa}

A relevância deste trabalho reside na demonstração prática de uma arquitetura de software altamente modular. O desenvolvimento de um algoritmo de xadrez do zero, com alta performance, diretamente em uma \textit{Game Engine} seria inviável para o escopo do projeto, além de computacionalmente ineficiente. Ao invés de reinventar ferramentas já consolidadas pela comunidade científica e de desenvolvedores, este projeto propõe uma ponte técnica elegante entre tecnologias de naturezas distintas (uma engine gráfica, uma linguagem de script utilitária e um executável externo). 

Essa abordagem arquitetural é fortemente justificável pois resolve uma lacuna comum no desenvolvimento de software de entretenimento: a dificuldade de executar processos pesados em segundo plano sem impactar a fluidez visual do sistema. Muitos desenvolvedores iniciantes enfrentam travamentos de tela ao tentar calcular lances de IA na mesma via de processamento (\textit{thread}) dos gráficos. Ao documentar a separação dessas responsabilidades, este projeto serve como um modelo de estudo valioso para estudantes de Análise e Desenvolvimento de Sistemas que precisem integrar motores visuais modernos a inteligências artificiais complexas.

\section{Contribuições realizadas}

O desenvolvimento deste simulador trouxe contribuições práticas e metodológicas para a estruturação de projetos de jogos de tabuleiro, destacando-se:

\begin{itemize}
	\item \textbf{Arquitetura de Comunicação Assíncrona:} A implementação de um script em Python que atua como ponte, traduzindo o estado do tabuleiro (formato FEN) e devolvendo a melhor jogada (formato UCI), provou que é viável acoplar programas externos pesados à Godot sem causar interrupções (\textit{stutters}) na interface principal.
	\item \textbf{Modularidade de Interface:} A centralização da navegação entre os diversos menus e telas do jogo por meio de um único gerenciador global demonstrou a aplicação de boas práticas de design de software, facilitando a expansão do código e evitando o acoplamento excessivo entre as telas.
	\item \textbf{Escalabilidade Visual:} A criação de um sistema de customização que utiliza estruturas de dados flexíveis (como Dicionários) substituiu lógicas condicionais repetitivas, resultando em um código mais limpo e otimizado para futuras atualizações de conteúdo estético.
\end{itemize}

\section{Organização do trabalho}

Este capítulo expôs o contexto inicial, o problema de pesquisa, os motivos, os objetivos e as contribuições deste trabalho. 
O capítulo \ref{capitulo:revisao_bibliografica} apresenta os principais conceitos teóricos sobre o xadrez computacional, assim como as ferramentas, linguagens e \textit{frameworks} utilizados para a solução do problema. 
O capítulo \ref{capitulo:materiais_metodos} apresenta a metodologia adotada e a arquitetura delineada para o desenvolvimento da aplicação. 
O capítulo \ref{capitulo:projeto} expõe o projeto de \textit{software} construído, detalhando os scripts, a organização visual e o fluxo de dados entre o jogo e a Inteligência Artificial. 
O capítulo \ref{capitulo:resultados_discussao} apresenta os resultados práticos obtidos com o simulador em funcionamento, discutindo a eficiência da comunicação assíncrona. 
Por fim, o capítulo \ref{capitulo:consideracoes_finais} compila as considerações finais sobre os desafios superados e sugere propostas para trabalhos futuros.
